\documentclass{plantillaPPT}
\titulo{11}{Convergencia Absoluta y Series de Potencias}
\begin{document}
\primeraDiapo

\begin{frame}{Criterio de Convergencia Absoluta}
\framesubtitle{Criterios de Convergencia}

Dada una sucesión \(a_n\) se tiene que:\\[2ex]
Si \(\displaystyle\sum_{n}^\infty |a_n|\) CV \(\Rightarrow\) \(\displaystyle\sum_{n}^\infty a_n\) CV\\[2ex]
Y se dice que la serie \(\displaystyle\sum_{n}^\infty a_n\) es \textbf{absolutamente convergente}.\\[4ex]
Además, si \(\displaystyle\sum_{n}^\infty a_n\) CV pero \(\displaystyle\sum_{n}^\infty |a_n|\) DV, se dice que la serie \(\displaystyle\sum_{n}^\infty a_n\) es \textbf{Condicionalmente convergente}.
\end{frame}

\begin{frame}{Criterio de la Razón/Cociente y de la Raíz}
\framesubtitle{Criterios de Convergencia}

Sea \(\displaystyle\sum_{n}^\infty a_n\) una serie de términos no nulos que cumple alguna de las expresiones:\\[2ex]

\[\lim_{n\to\infty}\frac{|a_{n+1}|}{|a_n|} = L \qquad \text{o bien} \qquad \lim_{n\to\infty}\sqrt[n]{|a_n|} = L\]\\[2ex]
Se tiene que:\\[2ex]
- Si \(L<1\) La serie Converge, y converge \textbf{absolutamente}. \\[2ex]
- Si \(L>1\) La serie Diverge.\\[2ex]
- Si \(L=1\) El criterio no es concluyente.
    
\end{frame}

\begin{frame}{Problema 1}
\framesubtitle{Práctica 11}

1. Determine si cada una de las siguientes series converge absolutamente, condicionalmente o no converge.

\begin{equation*}
    (a) \sum_{n=1}^\infty \frac{(n!)^2}{(2n)!} \qquad (b)\sum_{n=1}^\infty \frac{\ln(n)}{\sqrt[n]{n+1}} \qquad (c) \sum_{n=1}^\infty \frac{(-1)^nn}{e^{n^2}}
\end{equation*}

\end{frame}

\begin{frame}{Problema 2}
\framesubtitle{Práctica 11}

2. Determine, si es posible, los valores de $\alpha\in\mathbb{R}$ de modo que la serie
\begin{equation*}
    \sum_{n=1}^\infty \frac{\sqrt{n!}}{n^{\alpha n}}
\end{equation*}
sea convergente.

\end{frame}

\begin{frame}{Series de Potencia}
\framesubtitle{Definición}

\begin{definicion}
    Una \textbf{Serie de Potencia} centrada en $a$ es una serie de la forma:
    \[
    \sum_{n=0}^{\infty} c_n (x-a)^n = c_0 + c_1(x-a) + c_2(x-a)^2 + \dots
    \]
    donde $x$ es una variable y los $c_n$ son constantes llamadas coeficientes de la serie.
\end{definicion}

\textbf{Nota:} Si el centro es $a=0$, la serie se simplifica a $\sum_{n=0}^{\infty} c_n x^n$.
\end{frame}

\begin{frame}{Convergencia de Series de Potencia}
\framesubtitle{Radio e Intervalo de Convergencia}

\small % <--- Este comando reduce el tamaño de la fuente
\begin{enumerate}
    \item \textbf{Radio de Convergencia ($R$):} Es un número no negativo ($0 \le R \le \infty$) que determina el tamaño del dominio de convergencia. La serie:
    \begin{itemize}
        \item \textbf{Converge} absolutamente si $|x-a| < R$.
        \item \textbf{Diverge} si $|x-a| > R$.
    \end{itemize}
    \textcolor{blue}{Para Calcular $R$ se utiliza el Criterio de la Razón/Raíz.}
    \vspace{0.5cm}
    \item \textbf{Intervalo de Convergencia ($I$):} Es el conjunto de valores de $x$ para los cuales la serie converge. Siempre incluye el intervalo abierto $(a-R, a+R)$, y es \textbf{esencial} verificar la convergencia en los puntos extremos $x=a-R$ y $x=a+R$ por separado.
\end{enumerate}
\end{frame}

\begin{frame}{Problema 3}
\framesubtitle{Práctica 11}

3. Determine el radio de convergencia de las siguientes series de potencias
\begin{equation*}
\begin{gathered}
(a) \sum_{n=1}^\infty \frac{(-1)^n(x+1)^n}{n!5^n} \qquad (b) \sum_{n=2}^\infty \frac{(-1)^n\ln(2)(2x-1)^n}{n} \\[1cm]
(c) \sum_{n=1}^\infty \frac{(-3)^n x^n}{\sqrt{n+1}}
\end{gathered}
\end{equation*}

\end{frame}

\end{document}
